\chapter{Typed genomic computation systems}
Two arrays both have shape $2\times2$. In the first, rows are samples and
columns are features. In the second, rows are features and columns are
samples. A shape-only interface accepts one where it expects the other.
The matrix operation can run and return plausible numbers while answering
the wrong question.

This chapter builds a small language in which a computational candidate
declares its input meanings, operations, parameter owners and state effects.
We then execute a two-instruction genome, inspect its trace, and reject
malformed variants before numerical work. The word \emph{genomic} names
an Evolutor research abstraction here; these are software instructions,
not DNA molecules, biological genes or a trained genomic model.

Prerequisites are Chapter 12's graph construction and Chapter 13's
separation of structure, parameters and transient state. The deliverable is
an inspectable typed reference interpreter. Production compilation,
distributed effects, automatic differentiation and deployment are not
silently included.

\section{Start with a contract that shapes alone cannot express}
\index{semantic axis}
Define a tensor type as
\begin{equation}
 \tau=(A,S,d),
 \qquad A=(a_1,\ldots,a_r),\quad S=(s_1,\ldots,s_r),
 \label{eq:type}
\end{equation}
where $A$ lists semantic axis names, $S$ lists positive integer sizes,
$d$ is an element type and $r$ is rank. In this reference, names within a
type are distinct and $d=\texttt{float64}$. We allow rank-zero scalar types
but no zero-size or dynamic-size axes. Those are scope choices, not claims
that empty or dynamic tensors are invalid in general.

Type equality requires all three fields to agree in order. Thus a tensor
whose axes are \texttt{(sample, feature)} is not interchangeable with one
whose axes are \texttt{(feature, sample)}, even when both have sizes
$(2,2)$ and dtype \texttt{float64}.
An explicit transpose can convert both values and their types.
Relabeling metadata without transposing data is not that conversion.

\Plate{01-types}{Equal shape does not imply equal meaning. An explicit
transpose changes the array and swaps its semantic axes. The direct
connection is rejected by exact type matching even though a shape-only
runtime could execute it.}{fig:types}

\Listing{TensorType}

A NumPy array does not carry these semantic names. The runtime checks shape,
dtype and finiteness; the caller supplies the truthful semantic declaration.
The type system can check consistent use of that declaration, not prove
that a file labeled \texttt{feature} actually contains the intended features.
This trust boundary belongs beside the type definition.

\section{Define the alphabet and the computational gene}
\index{identity bypass}
Our operation alphabet is deliberately small:
\begin{equation}
 \Sigma=\{\texttt{scale},\texttt{accumulate}\}.
\end{equation}
A \Term{computational gene} is an immutable instruction record
$g=(\mathrm{id},\mathrm{op},\mathrm{src},\mathrm{dst},\tau,\mathrm{owner})$.
The record names a source register and a fresh destination, declares their
common type, and identifies an external parameter or state owner.
A \Term{genome} is an ordered tuple $(g_1,\ldots,g_L)$.
It is not yet an unordered graph or a self-modifying program.

\Listing{Gene}

For \texttt{scale}, the owner names a finite scalar parameter $p$:
$y=px$. For \texttt{accumulate}, it names a state tensor $h$ of type $\tau$:
$y=x+h$ and a proposed next state $h'=y$. State has no trainable update rule
in this chapter. Parameter values are read-only during expression.

\Plate{02-genome}{A genome owns an ordered collection of immutable
instruction records. Genes reference external parameter or state owners;
they do not duplicate those objects. Regulator decisions select whether
each instruction uses its active operation or its declared bypass.}{fig:genome}

Separate \Term{regulation} from instruction validity. Our regulator provides
one explicit Boolean decision per gene per expression step. When disabled,
either operation returns a copy of its input; an accumulating gene also
leaves its state unchanged. This identity bypass is part of the language
contract. Other languages might omit an output or return zeros, but cannot
leave that choice implicit. A disabled malformed gene is still rejected:
regulation is not a way to hide an invalid instruction.

The example genome contains
\[
 g_1:\ u=\texttt{scale}_{\mathrm{gain}}(x),\qquad
 g_2:\ y=\texttt{accumulate}_{\mathrm{memory}}(u).
\]
Here $x,u,y,h\in\mathbb R^2$ share axis \texttt{feature},
$\mathrm{gain}=2$, and initial memory is $(0,0)$.

\section{Type checking as an environment transformation}
\index{type preservation}\index{state ownership}
Let $\Gamma$ map defined value names to their tensor types, $\Theta$ contain
declared parameter owners, and $\Delta$ map state owners to tensor types.
A scale instruction is well-formed when
\begin{equation}
 \Gamma(\mathrm{src})=\tau,\quad
 \mathrm{dst}\notin\operatorname{dom}\Gamma,\quad
 \mathrm{owner}\in\Theta.
\end{equation}
It extends the environment with
$\Gamma'=\Gamma[\mathrm{dst}\mapsto\tau]$.
An accumulation additionally requires
$\Delta(\mathrm{owner})=\tau$ and no earlier instruction writing that state
owner in this expression. Gene identities must also be unique.

These are \Term{single static assignment} (SSA) value names: a value name
is defined once. State ownership is a different constraint. Many genes
may read the same scalar parameter; this reference permits only one
writer to each state owner. A more expressive language could order multiple
writes explicitly, but merely listing them in a file would not specify
whether each reads the old or newly updated state.

\Listing{check_genome}
\Plate{04-rejection}{The type environment grows only after an instruction
passes its checks. Missing producers, incompatible semantic axes, reused
destinations and undeclared owners fail before arithmetic. State-writer
conflicts are checked independently of value-name uniqueness.}{fig:rejection}

The checking pass is linear in instruction count under expected constant-time
dictionary operations, apart from type comparisons proportional to rank.
This is a logical work estimate, not a compiler benchmark.

A small preservation argument explains why the checks matter. Assume all
inputs satisfy their declared types. A finite scalar multiplication preserves
shape, axis interpretation and element dtype in this reference; addition of
equal-typed arrays does likewise. Bypass copies the input. Each successful
instruction therefore produces a value of its declared type. By induction,
all defined values remain well-typed. Floating-point overflow can still
violate finiteness, so the runtime checks each result and fails rather than
committing invalid state. Type correctness is not numerical stability.

\section{Expression has a snapshot and a commit boundary}
\index{snapshot semantics}\index{aliasing}
For one expression step, the environment consists of input values,
parameter values, a state snapshot and regulator decisions. Evaluation
reads only the old snapshot $h_t$ while producing a proposed $h_{t+1}$.
The caller adopts the returned state only after successful completion.
The input arrays and caller-owned state arrays are never mutated.

\Plate{03-step}{The two-instruction example reads a stable state snapshot.
With input $(1,2)$ and gain 2, the first step proposes state $(2,4)$ and
the second proposes $(4,8)$. Disabling accumulation on the third step
bypasses to $(2,4)$ while retaining memory $(4,8)$.}{fig:step}

The public \texttt{express} function checks the complete genome, exact
input/state key sets, all runtime values, scalar parameters and an explicit
Boolean decision for every gene. It then calls this internal evaluator.
Calling the helper directly bypasses validation and is not a supported
untrusted-input entry point.

\Listing{evaluate_checked}
\begin{center}\input{\Generated/step-table.tex}\end{center}

No mutation is committed if an exception occurs midway. This is an
\Term{atomic state transition} at the reference API boundary, not a durable
database transaction. It offers no process-crash recovery, cross-worker
consensus or protection against concurrent mutation by the caller.
The snapshot contract assumes no external thread changes the arrays or
parameters during a call.

Why copy output into next state? If the returned output and returned state
shared writable storage, a caller changing its output buffer could change
future model behavior. The reference tests that mutating a returned value
does not mutate the returned state or original input.

Why read \texttt{states}, not \texttt{next\_state}? This makes snapshot
semantics explicit. Our one-writer restriction keeps the small language
simple, but the distinction remains essential when later languages
introduce more complex read/write dependencies.

\section{A trace is evidence only if it identifies the experiment}
\index{trace!replay}
The reference emits a compact per-gene trace containing the supplied genome
version, gene identity, enabled decision and numeric output. This is enough
to inspect our small worked example, not enough to replay an arbitrary
experiment without its external context.

\Plate{05-trace}{Replay needs the genome, inputs, parameter values,
initial state, regulator decisions and execution rules. The compact
reference trace records only a subset of this context. A version string
is a label, not a content hash or tamper-evident signature.}{fig:trace}

A practical replay record should bind immutable content identities for
the genome, type declarations, parameter snapshot, input batch, initial
state and regulator policy or realized decisions. Also record numeric
precision, operator implementations, random seeds and environment.
For stochastic operations, a seed alone may be insufficient if kernel
scheduling changes random-number consumption.

Reproducibility has levels. A semantic replay can agree within a tolerance;
a bitwise replay must agree exactly. Neither proves biological usefulness.
A trace can establish that a specified program ran on specified data, while
an external evaluation must establish whether its answers are good.

\section{Current compiler practice and the research boundary}
\index{compiler!legality}
The current StableHLO specification defines programs using typed functions
and operations, with tensor shapes and element types; it also distinguishes
token values used for ordering \cite{stablehlo}. That separation is relevant
to the distinction between data and effects here. Our named semantic axes,
gene records and snapshot policy are original teaching choices, not
StableHLO syntax or a compatible dialect.

The industry-facing lesson is to make interfaces machine-checkable before
attempting optimization. A lowering must preserve not only array shapes but
also state-update order, parameter sharing, bypass semantics and failure
behavior. Fusing two operations can be invalid if the unfused program rejects
a nonfinite intermediate while the fused program hides it. Reordering
stateful operations can change results even when their arithmetic costs match.

Evolutor remains the broader research/compiler/runtime direction above
DOGMA (non-Transformer DNA-native model plus DOGMA Engine) and Hermon DNA
(Transformer-based DNA model plus Hermon DNA Engine). This chapter supplies
neither engine. It supplies a small contract against which future compiler
transformations and runtime implementations can be tested.

\section{Exercises with worked answers}
\begin{enumerate}
\item \textbf{Type mismatch.} Can two $2\times2$ arrays be incompatible?
\textit{Answer:} Yes. Reversing \texttt{sample} and \texttt{feature} changes
the declared meaning despite equal dimensions.
\item \textbf{Conversion.} Is changing axis labels sufficient to transpose?
\textit{Answer:} No. Transpose changes the coordinate map of the data as
well as the type. Metadata-only relabeling can assert a false interpretation.
\item \textbf{Environment.} Why must a destination be fresh?
\textit{Answer:} SSA avoids ambiguous value versions. Reusing \texttt{x}
could conceal whether later instructions read the original or overwritten value.
\item \textbf{State.} Compute two steps with input $(1,2)$, gain 2 and
initial state zero. \textit{Answer:} Each scale gives $(2,4)$.
Accumulation yields states $(2,4)$ then $(4,8)$.
\item \textbf{Bypass.} Disable accumulation on the next step.
\textit{Answer:} Output is the scaled input $(2,4)$; memory remains $(4,8)$.
Disabled does not mean output zero or memory reset.
\item \textbf{Ownership.} Can two scale genes share one parameter owner?
\textit{Answer:} Yes; both read the same scalar. Two accumulating genes
cannot share one state owner under this reference's one-writer rule.
\item \textbf{Negative control.} A disabled gene has an undeclared owner.
Should it pass? \textit{Answer:} No. Genome validity is checked independently
of today's regulator decisions.
\item \textbf{Numerics.} Can a well-typed scale fail at runtime?
\textit{Answer:} Yes. Finite inputs and a finite parameter can overflow
their product. The finite-result check rejects it without mutating caller state.
\item \textbf{Aliasing.} Why copy the accumulated output into next state?
\textit{Answer:} Otherwise a later caller mutation of the output array could
alter stored memory through a shared buffer.
\item \textbf{Replay.} Is the version string alone sufficient?
\textit{Answer:} No. Bind content, values, decisions and execution rules.
A mutable label identifies no immutable artifact.
\item \textbf{Extension.} Add an explicit transpose instruction.
\textit{Rubric:} Validate a rank permutation; derive output axes and sizes;
transpose the value consistently; reject duplicate/out-of-range axes;
test inverse permutation and a same-shape semantic mismatch.
\item \textbf{Compiler test.} Propose a legality test for fusion.
\textit{Rubric:} Compare outputs, proposed state, bypass decisions,
parameter sharing and failures against the reference. Include nonfinite
intermediates, disabled genes and rejected ownership conflicts.
State the numerical tolerance and what it cannot prove.
\end{enumerate}

\section{The next step is an operational semantics}
\textbf{Reproduce the chapter.} From the repository root, run
\texttt{python3 drafts/ch14/build.py --assets-only}, then
\texttt{python3 -m pytest tests/test\_ch14\_reference.py}.
The first command regenerates \texttt{results.json}, tables, plots and literal
code listings; the tests compare the implementation with independent
oracles and explicit failure cases. Run the builder without
\texttt{--assets-only} to compile this review PDF.

The alphabet, instruction records, type environment, regulator decisions,
state snapshot and trace now have distinct meanings. This prevents a
biological analogy from standing in for an executable specification.

\textbf{Working vocabulary.} A tensor type declares axes, sizes and dtype;
a computational gene names an operation and its dependencies; a genome
orders genes; expression executes one regulated step; SSA versions values;
ownership identifies shared resources; atomic return separates proposed
from adopted state. Chapter 15 will turn these contracts into explicit
transition rules and expression traces, then ask which different execution
orders are genuinely equivalent.
